
\documentclass[11pt]{article}

\usepackage[margin=1in]{geometry}
\usepackage{amsmath, amssymb, amsthm}
\usepackage{mathtools}
\usepackage{enumitem}
\usepackage{hyperref}
\usepackage{microtype}

\hypersetup{
  colorlinks=true,
  linkcolor=blue,
  citecolor=blue,
  urlcolor=blue
}

\title{From Rewrite Semantics to Evidence:\\
A Conceptual Primer Linking GSLT, OSLF, Behavioral Types,\\
World-Model Evidence, and Knuth--Skilling Measurement\\
\large (Conversation Primer, 2026-02-14)}
\author{}
\date{}

\newcommand{\Psh}[1]{\widehat{#1}} % presheaf category notation (mnemonic)
\newcommand{\Sub}{\mathrm{Sub}}
\newcommand{\Hom}{\mathrm{Hom}}
\newcommand{\Obj}{\mathrm{Obj}}
\newcommand{\id}{\mathrm{id}}

\newcommand{\DiamondOp}{\mathbin{\Diamond}}
\newcommand{\BoxOp}{\mathbin{\Box}}
\newcommand{\forces}{\models}

\newcommand{\Red}{\;\Rightarrow\;}
\newcommand{\RedStar}{\;\Rightarrow^{*}\;}
\newcommand{\entails}{\vdash}
\newcommand{\meet}{\wedge}
\newcommand{\join}{\vee}
\newcommand{\botEl}{\bot}
\newcommand{\topEl}{\top}

\theoremstyle{definition}
\newtheorem{definition}{Definition}
\newtheorem{remark}{Remark}
\newtheorem{principle}{Principle}
\newtheorem{goal}{Target Theorem}

\begin{document}
\maketitle

\vspace{-1.0em}

\begin{abstract}
This note distills a conceptual architecture for connecting (i) operational/rewrite semantics of languages,
(ii) an \emph{intrinsic} logic of behavioral properties and ``native types'' derived from that semantics,
(iii) an observer- and evidence-aware world-model semantics, and (iv) quantitative measurement
principles for plausibility/evidence that justify when scalar probabilities are appropriate and when richer
(partial-order, multi-dimensional, or imprecise) valuations are structurally necessary.
The aim is to serve as a didactic primer for future discussions and formalization work,
while avoiding overclaims: each layer has its own truth conditions, computational constraints,
and appropriate ``claim hygiene.''
\paragraph{Status note.} Implementation-status and count claims in this draft are snapshot-scoped; verify against current repository trackers and build outputs before citing as current state. Lean-proven claims are only those tied to explicit Lean file/theorem references; broader mathematical or historical claims are literature-backed context and not asserted as Lean-proven unless explicitly stated.
\end{abstract}

\tableofcontents

\section{Executive summary: the three (plus one) layers}

The unifying picture has three main layers (plus an optional fourth):

\begin{enumerate}[leftmargin=1.5em]
\item \textbf{Behavior layer (operational semantics).}
A language is presented as a rewrite/reduction system.
States are terms/configurations; steps are reductions.
This is where ``what can happen'' is defined.

\item \textbf{Behavioral logic/type layer (OSLF-style).}
From a reduction relation, one constructs a logical space of predicates on states
(often organized as a Heyting algebra), together with modal operators induced by the reduction structure.
This is a \emph{behavioral} semantics: it talks about reachability, invariants, and interaction structure,
not about denotations in a separate model.

\item \textbf{Measurement layer (world-model evidence + valuation theory).}
A world model assigns \emph{evidence} to queries about behavior.
Evidence lives in a structured algebra (often richer than $\mathbb{R}$), supporting revision/combination.
Quantitative plausibility constraints (Knuth--Skilling style) explain when and why
regraduation to a scalar probability-like scale is justified, and when it is not.

\item \textbf{(Optional) Observer/indistinguishability layer (distinctions).}
One may explicitly model an observer's (in)ability to distinguish states, yielding
graphs or relations of indistinguishability and entropy-like summaries of the observer's discrimination.
This layer is useful for resource-bounded agents and imprecise evidence.
\end{enumerate}

A core methodological stance is:
\begin{quote}
\emph{Keep the behavioral logic honest and compositional;
keep the measurement layer explicit about assumptions;
avoid ``smuggling'' an oracle or a scalar collapse into the semantics without stating it.}
\end{quote}

\section{GSLT-style language specifications: rewrite systems with premises}

\subsection{Why a rewrite-first presentation?}

Many systems of interest---programming languages, process calculi, grammar formalisms,
inference engines---admit a presentation where the \emph{dynamics} is given by rewriting:
\begin{itemize}[leftmargin=1.5em]
\item a syntax of terms or configurations (often multi-sorted),
\item equations / structural congruence (rearrangements that preserve meaning),
\item rewrite rules (directed dynamics),
\item \emph{premises} (side conditions) such as freshness, matching constraints,
or oracle relations (e.g., a grounded/native function call).
\end{itemize}

This style is attractive because it provides a single substrate for:
\begin{enumerate}[leftmargin=1.5em]
\item an executable engine (compute next steps),
\item a declarative relation (what \emph{counts} as a step),
\item and higher-level constructions (modalities, type systems, evidence semantics).
\end{enumerate}

\subsection{Premise-aware rewriting is not optional}

Realistic specifications frequently require side conditions:
fresh-name generation, matching constraints, or external computations.
A key distinction is:
\begin{itemize}[leftmargin=1.5em]
\item \textbf{Internal premises:} decidable properties computed on syntax (freshness, equality, matchability).
\item \textbf{Oracle premises:} relations answered by an external environment
(grounded calls, database queries, probabilistic sensors).
\end{itemize}

A premise-aware rewrite semantics should be defined in \emph{two} aligned ways:
\begin{enumerate}[leftmargin=1.5em]
\item a \textbf{declarative} step relation \,$p \Red q$\, that explicitly quantifies over premise witnesses;
\item an \textbf{executable} step function that computes candidate reducts and is proved
sound (every computed reduct is valid) and complete (every valid reduct is computed),
relative to the premise environment.
\end{enumerate}

This ``engine $\leftrightarrow$ declarative'' equivalence becomes a central lemma reused throughout the stack.

\section{OSLF: behavioral logic from operational semantics}

\subsection{The minimal construction: modalities from a reduction relation}

Let $S$ be a set (or type) of states and let $R \subseteq S \times S$ be a one-step reduction relation.
A \emph{process predicate} is a subset $\varphi \subseteq S$ (or a function $S \to \mathsf{Prop}$).

Define modal operators on predicates by:
\begin{align*}
(\DiamondOp \varphi)(p) &:\iff \exists q,\; R(p,q) \;\wedge\; \varphi(q),\\
(\BoxOp \psi)(p) &:\iff \forall q,\; R(p,q) \to \psi(q).
\end{align*}

\begin{principle}[Galois connection]
Under these definitions, $\DiamondOp$ and $\BoxOp$ form a Galois connection:
\[
\DiamondOp \varphi \subseteq \psi \quad\text{iff}\quad \varphi \subseteq \BoxOp \psi.
\]
\end{principle}

This single adjunction is the seed of a large amount of structure:
it yields monotonicity, distributive properties, and a disciplined language of reachability
and invariance claims.

\subsection{From predicates to a Heyting algebra of ``types''}

Predicates ordered by implication/inclusion form a complete lattice (under mild assumptions).
With implication interpreted as Heyting implication, one obtains a Heyting algebra:
\begin{itemize}[leftmargin=1.5em]
\item $\varphi \meet \psi$ = conjunction (intersection),
\item $\varphi \join \psi$ = disjunction (union),
\item $\varphi \Rightarrow \psi$ = Heyting implication (weakest predicate making $\varphi$ entail $\psi$),
\item $\topEl$ and $\botEl$ = tautology and contradiction.
\end{itemize}

\begin{remark}
This is a \emph{logic of behavior}: propositions are predicates on states.
It is not, by itself, a denotational model assigning external meanings to programs.
\end{remark}

\subsection{The topos upgrade: subobject classifiers and presheaf semantics}

The ``predicate = subset'' picture is Set-level.
A more categorical presentation treats predicates as \emph{subobjects} and organizes them via the \emph{subobject fibration}.
For presheaf categories $\Psh{C} \simeq C^{op} \to \mathsf{Type}$, predicates correspond to subpresheaves,
and the subobject classifier $\Omega$ is described by sieves.

The conceptual steps (useful as a formalization roadmap) are:
\begin{enumerate}[leftmargin=1.5em]
\item construct $\Omega$ in $\Psh{C}$ (sieves object),
\item show the subobject presheaf $\Sub(-)$ is representable by $\Omega$,
\item obtain \textbf{HasClassifier} for the presheaf topos constructively,
\item define change-of-base along morphisms and derive $\exists_f \dashv f^* \dashv \forall_f$,
\item prove Beck--Chevalley conditions for pullback squares,
\item instantiate the OSLF ``modalities from a span'' story internally in the topos/hyperdoctrine.
\end{enumerate}

\begin{remark}
In practice, it is helpful to keep a clear separation between:
(i) the constructive presheaf/topos proofs (classifier, adjoints, Beck--Chevalley),
and (ii) the operational reduction engine and its soundness/completeness.
The full story requires both.
\end{remark}

\section{Behavioral ontology: bisimulation, logical equivalence, and hypotheses}

\subsection{Bisimulation as ``semantic indistinguishability''}

Given a transition system $(S,R)$, a (strong) bisimulation is a relation $\sim$ such that
$p \sim q$ implies matching step structure in both directions.
Bisimulation is typically an equivalence relation, enabling the quotient $S/{\sim}$.
Intuitively:
\begin{quote}
\emph{Bisimulation classes are candidates for ``behavioral individuals.''}
\end{quote}

\subsection{Modal logics and characterization theorems}

Modal formulas built from $\DiamondOp$, $\BoxOp$, and connectives are usually \emph{bisimulation-invariant}:
if $p \sim q$ then $p$ and $q$ satisfy the same modal formulas.
A deeper question is \emph{completeness}:
does ``satisfies the same formulas'' imply bisimulation?

A classical phenomenon (often associated with Hennessy--Milner style results) is:
\begin{itemize}[leftmargin=1.5em]
\item Under \textbf{image-finiteness} (finite branching) or similar hypotheses,
a finitary modal logic can characterize bisimulation.
\item Without such hypotheses, one often needs \textbf{infinitary} conjunction/disjunction
and/or \textbf{$\mu$-calculus} fixpoints to fully characterize bisimulation.
\end{itemize}

\begin{remark}[Philosophical reading]
Image-finiteness can be read as a \emph{bounded-choice} assumption:
finite tests can in principle separate behaviors because there is not an infinitary fan-out at each step.
When branching is unbounded or behaviors are inherently recursive/reflective,
finitary observation is not complete without enrichment (infinitary or fixpoint structure).
\end{remark}

\section{World-model semantics: evidence, revision, and guarded rewrites}

\subsection{Why a world model?}

OSLF gives a compositional logic of behavioral predicates.
But an agent typically needs more than boolean truth:
it needs \emph{graded evidence} under uncertainty, partial observability, and correlation.

A world model provides:
\begin{itemize}[leftmargin=1.5em]
\item a state space $W$ (often containing a joint distribution or structured uncertainty),
\item a query language $Q$ (questions about the world or computation),
\item an evidence structure $E$,
\item an evidence extraction map $\mathsf{evidence} : W \times Q \to E$,
\item a revision/combination operation on world states or evidences.
\end{itemize}

\subsection{Evidence quantales and Heyting-valued semantics}

A common pattern is to let $E$ be a (commutative) quantale or related algebra:
a partially ordered space with combination and joins,
supporting both aggregation and revision.
In Heyting-compatible settings, connectives are interpreted by the lattice/Heyting operations.

\begin{principle}[No premature scalar collapse]
If evidence is naturally partial-order valued (multi-dimensional, interval/imprecise, or structured),
collapsing to a single real number without stated assumptions generally loses information.
Scalar probabilities arise as a \emph{view} under additional constraints.
\end{principle}

\subsection{Guarded rewrite calculus for queries}

A practical and honest approach is to use a rewrite calculus for queries with explicit side conditions $\Sigma$:
\begin{quote}
\emph{Rules are applied only when stated assumptions about independence, approximation, or model class hold.}
\end{quote}
This enables:
\begin{itemize}[leftmargin=1.5em]
\item derivations that are sound relative to the world model and side conditions,
\item explicit tracking of approximation,
\item modular composition of inference heuristics with a small trusted kernel.
\end{itemize}

\section{Knuth--Skilling measurement: when do numbers appear?}

\subsection{Two complementary roles of KS-style theory}

Knuth--Skilling style axiomatizations are useful in two ways:
\begin{enumerate}[leftmargin=1.5em]
\item \textbf{Measurement value theory:} if plausibility values have certain order/combination symmetries,
there exists a regraduation into a real scale where combination becomes additive (or otherwise tractable).
\item \textbf{Valuation on event/proposition lattices:} if events form a (distributive/Boolean) lattice,
then plausible axioms (sum rules, chain rules) yield probability-like calculus, including Bayes-like laws under suitable structure.
\end{enumerate}

\subsection{Totality vs.\ imprecision}

A recurring technical and philosophical distinction is:
\begin{itemize}[leftmargin=1.5em]
\item \textbf{Total/linear order assumptions} enable faithful embedding into $\mathbb{R}$
and classical probability-like behavior.
\item \textbf{Partial orders} naturally support imprecise or multi-dimensional evidence;
representation by a single scalar is generally not faithful without loss.
\end{itemize}

In the integrated picture, one uses KS-style results to justify \emph{which} scalar views are legitimate,
and to document the assumptions needed for each view.

\section{Distinction graphs: observer-relative indistinguishability and graphtropy}

\subsection{From equivalence relations to graphs}

If indistinguishability is transitive, it forms an equivalence relation and yields a partition.
If indistinguishability is not transitive (as with resource-bounded or noisy observers),
it is naturally represented as a graph of non-distinction, possibly weighted.
This complements bisimulation:
\begin{itemize}[leftmargin=1.5em]
\item \textbf{Bisimulation} is an ``ideal observer'' equivalence induced by the semantics.
\item \textbf{Distinction graphs} model an actual observer's discrimination regime (tests, thresholds, noise).
\end{itemize}

\subsection{Graphtropy as measurement of discrimination}

Graph-based generalizations of logical entropy (``graphtropy'') summarize how much an observer distinguishes,
and weighted versions can incorporate graded distinguishability.
This provides a bridge between qualitative behavioral ontology and quantitative measurement,
especially when scalar probability is not the right primitive.

\section{Integration blueprint: a clean end-to-end story}

\subsection{A canonical pipeline}

A disciplined, non-overclaimed pipeline looks like:

\begin{enumerate}[leftmargin=1.5em]
\item \textbf{Language spec (GSLT)}: syntax + equations + rewrite rules with premises.
\item \textbf{Operational engine}: executable one-step search (premise-aware), with soundness/completeness to a declarative step relation.
\item \textbf{Behavioral logic (OSLF)}:
  \begin{itemize}[leftmargin=1.5em]
  \item predicate space (Heyting structure),
  \item modalities induced by reduction (Galois connection),
  \item (optional) topos/hyperdoctrine internalization.
  \end{itemize}
\item \textbf{World model}: evidence-valued semantics for atoms/queries; revision rules.
\item \textbf{Measurement constraints}: KS-style regraduation/valuation theorems justify scalar views and clarify when imprecision is required.
\item \textbf{Tools}: bounded checkers, abstract interpreters, and executable analyzers as approximations with stated correctness guarantees.
\end{enumerate}

\subsection{Two-track R\&D that meets in the middle}

Two complementary development tracks help avoid both ``math-only'' and ``tool-only'' traps:

\paragraph{Track A: principled math-first}
Prioritize constructive topos/hyperdoctrine foundations:
classifier, adjoints, Beck--Chevalley, and internal change-of-base.
This ensures the derived modalities and quantifiers are not merely asserted but constructed.

\paragraph{Track B: executable checker-first}
Prioritize premise-aware rewriting engines, bounded checkers, and agreement tests.
Prove soundness (and completeness where feasible) so the executable tool is not a black box.

\paragraph{Where they meet}
The meeting point is an end-to-end theorem: ``checker says \textsf{sat} $\Rightarrow$ semantic satisfaction,''
followed by transport into world-model evidence semantics (possibly via thresholding or an evidence-to-Prop view).

\section{MeTTa and GF as motivating clients}

\subsection{MeTTa as a GSLT: state machines + oracle boundaries}

A key modeling lesson is to treat a realistic language like MeTTa as a \emph{state transition system}:
\begin{itemize}[leftmargin=1.5em]
\item states include plan/stack, atomspace, bindings, and intermediate results;
\item rewriting encodes one machine step;
\item grounded/native calls are handled as explicit oracle premises.
\end{itemize}
This preserves honesty: the semantics does not pretend to compute what must be queried.

\subsection{GF as a query/meaning interface}

Grammar formalisms can be treated as rewrite/derivation systems producing structured meanings and paraphrases.
This creates opportunities to:
\begin{itemize}[leftmargin=1.5em]
\item express semantic queries as OSLF formulas over derivation dynamics,
\item interpret queries evidentially in a world model,
\item compare linguistic variants by behavioral invariants (language-independent abstract structures).
\end{itemize}

\section{Claim hygiene: avoiding facepalmable statements}

\begin{principle}[Do not conflate layers]
Avoid claims of the form ``we derived X'' when what is implemented is:
\begin{itemize}[leftmargin=1.5em]
\item hand-defined objects + proofs they satisfy an abstract interface, rather than
\item an executable algorithm producing those objects from the language spec.
\end{itemize}
Both are valuable, but they are different claims.
\end{principle}

\begin{principle}[Make oracle boundaries explicit]
If grounded calls, external databases, or probabilistic sensors are involved,
keep them as explicit premises/oracles with a clearly stated interface and semantics.
\end{principle}

\begin{principle}[Document assumptions for scalar views]
If a scalar probability/plausibility scale is used, document the order/separation/regularity assumptions
and present scalarization as a \emph{view} on richer evidence unless proven otherwise.
\end{principle}

\section{A minimal ``next theorem chain'' to target}

The following chain is a compact target for a rigorous end-to-end bridge:

\begin{goal}[Engine--declarative equivalence]
Premise-aware executable reduction is sound and complete for the declarative relation.
\end{goal}

\begin{goal}[OSLF adjunction]
Modal operators induced from the reduction relation satisfy $\DiamondOp \dashv \BoxOp$.
\end{goal}

\begin{goal}[Topos foundations (constructive)]
In the presheaf topos: $\Omega$ exists, $\Sub(-)$ is representable,
change-of-base has $\exists \dashv f^* \dashv \forall$, and Beck--Chevalley holds for pullbacks.
\end{goal}

\begin{goal}[Checker soundness]
If the bounded checker returns \textsf{sat}, then semantic satisfaction holds.
\end{goal}

\begin{goal}[Evidence semantics bridge]
Define an evidence-valued interpretation of atoms and lift it to formulas.
Show checker soundness implies a thresholded (or otherwise structured) evidence judgment.
\end{goal}

\begin{goal}[Measurement justification]
Under explicit KS-style assumptions, prove a regraduation/representation theorem
showing when scalar probability-like semantics is a conservative view.
\end{goal}

\section*{Acknowledgements and scope}
This note is intentionally conceptual and claim-hygienic.
It is not a substitute for the detailed formal developments in a proof assistant,
nor does it assert that every item above is already implemented.
Its purpose is to provide a shared map for future formalization and discussion.

\section*{Selected references (non-exhaustive)}
\begin{itemize}[leftmargin=1.5em]
\item S.~Mac Lane and I.~Moerdijk, \emph{Sheaves in Geometry and Logic}.
\item P.~T.~Johnstone, \emph{Sketches of an Elephant: A Topos Theory Compendium}.
\item B.~Jacobs, \emph{Categorical Logic and Type Theory}.
\item M.~Hennessy and R.~Milner, work on modal logic and bisimulation (Hennessy--Milner theorem phenomenon).
\item Knuth--Skilling style work on measurement foundations of inference (regraduation, valuation theory).
\item Work on subobject classifiers and presheaf toposes; see also formalization libraries (mathlib) for constructive developments.
\item Work on logical entropy (Ellerman) and observer-relative distinction frameworks (distinction graphs, graphtropy).
\end{itemize}

\end{document}
